\chapter{Application security and DevSecOps integration}

\section{Engineering principles}

Use agents to contextualize scanner findings, not replace scanners. SAST examines code patterns and flows; DAST observes a running target; SCA identifies dependency exposure; secret scanning detects credential candidates; infrastructure and container scans examine configuration and images. Each produces evidence with false positives and coverage limitations. A remediation agent must respect the finding scope and avoid exposing sensitive material.\src{W1,G2}

\section{Implementation approach}

Normalize findings by scanner, rule, affected location, commit, severity and evidence reference. Ask the agent to inspect reachability, propose a fix and add a regression test in an isolated environment. Re-run the same scanner and verify the original vulnerable behavior is absent. For a suspected secret, redact the value, notify the owner and use a separate approved revocation process; deletion from Git does not revoke it.

\section{Verification and limitations}

Track false-positive dispositions with rationale and reviewer identity. Do not authorize an agent to suppress rules or lower severity merely to make checks pass. Restrict DAST to explicitly authorized test targets.

\section{Chapter conclusion}

Apply the stated verification criteria before expanding the scope of this workflow. Record deviations, evidence and unresolved risks in the task record.
